\documentclass[11pt, a4paper]{article}

\usepackage[a4paper, margin=2.5cm]{geometry}
\usepackage{amsmath, amssymb}
\usepackage{hyperref}

\title{Topological Integrity as a Global Constraint in Gravitational Theory}
\author{Kai Stefan Dietrich \\ Independent Research}
\date{\today}

\begin{document}

\maketitle

\begin{abstract}
We propose a spectral-geometric framework in which physically admissible spacetime configurations 
are constrained by a global condition of topological integrity. In addition to curvature-squared 
terms, we include gradient-based contributions that encode variations of curvature. The resulting 
field equations extend General Relativity by incorporating both amplitude and variation of curvature. 
We derive explicit cosmological equations in the spatially flat FLRW case and discuss implications 
for black hole structure, cosmology, and gravitational stability.
\end{abstract}

% --------------------------------------------------
\section{Introduction}

General Relativity provides a geometrical description of gravitation 
in terms of spacetime curvature \cite{einstein1915, wald1984}. 

Black hole thermodynamics suggests deeper links between geometry 
and information \cite{bekenstein1973, hawking1975}. 

Cosmological models remain dependent on parameterized inputs 
rather than structural necessity \cite{friedmann1922}. 

Quadratic curvature corrections have also appeared in modified-gravity contexts \cite{starobinsky1980}. 

These observations motivate the introduction of a global consistency principle 
in which admissible configurations are selected by geometric constraints.

% --------------------------------------------------
\section{Topological Integrity Formalism}

We define an integrity functional:

\begin{equation}
\mathcal{I}[g] = \int_M \mathcal{F}\, d\mu_g
\end{equation}

with a minimal local realization:

\begin{equation}
\mathcal{F} = \alpha R^2 + \beta \nabla_\mu R \nabla^\mu R
\end{equation}

where $\alpha$ and $\beta$ are coupling parameters.

% --------------------------------------------------
\section{Field Equations}

We consider the action:

\begin{equation}
S =
\int d^4x \sqrt{-g}
\left[
\frac{1}{16\pi}(R - 2\Lambda)
+
\alpha R^2
+
\beta \nabla_\mu R \nabla^\mu R
\right]
+
S_{\mathrm{matter}}
\end{equation}

Variation yields:

\begin{equation}
G_{\mu\nu}
+
\Lambda g_{\mu\nu}
+
16\pi \alpha H_{\mu\nu}
+
16\pi \beta K_{\mu\nu}
=
8\pi T_{\mu\nu}
\end{equation}

with:

\begin{equation}
H_{\mu\nu}
=
2 R R_{\mu\nu}
-
\frac{1}{2} g_{\mu\nu} R^2
-
2 \nabla_\mu \nabla_\nu R
+
2 g_{\mu\nu} \Box R
\end{equation}

\begin{equation}
K_{\mu\nu}
=
\nabla_\mu R \nabla_\nu R
-
\frac{1}{2} g_{\mu\nu} (\nabla R)^2
-
\nabla_\mu \nabla_\nu R
+
g_{\mu\nu} \Box R
\end{equation}

% --------------------------------------------------
\section{Entropy and Geometric Integrity}

The $R^2$ term suppresses large curvature amplitudes, while the gradient term penalizes rapid variations of curvature. Together, these contributions define a balance between geometric concentration and dispersion.

% --------------------------------------------------
\section{Black Hole Solutions}

For Schwarzschild vacuum:

\begin{equation}
R = 0
\end{equation}

Thus corrections vanish externally, preserving classical solutions.

In high-curvature regimes, corrections become relevant and may regularize singularities.

% --------------------------------------------------
\section{Cosmology and Friedmann Dynamics}

For flat FLRW:

\begin{equation}
R = 6(\dot{H}+2H^2)
\end{equation}

The modified Friedmann equation becomes:

\begin{equation}
H^2 =
\frac{
8\pi \rho + \Lambda + 8\pi\alpha R^2 - 48\pi\alpha H \dot{R} + \beta (\dot{R})^2
}{
3(1 + 32\pi\alpha R)
}
\end{equation}

% --------------------------------------------------
\section{Predictions}

\begin{itemize}
\item Stability of spacetime under curvature fluctuations
\item Suppression of singularities
\item Modified cosmological evolution
\item Potential deviations in gravitational waves
\end{itemize}

% --------------------------------------------------
\section{Conclusion}

We presented a modified gravitational framework combining curvature amplitude and variation terms. The resulting theory remains consistent with General Relativity in appropriate limits while introducing new dynamical effects in high-curvature regimes.

% --------------------------------------------------
\bibliographystyle{unsrt}
\bibliography{references}

\end{document}
